% Problema 4, parte 1: campi da incollare nella piattaforma (uno per sezione)

% ===== CAMPO: 1. Risultato e ambito =====
Tre classi disgiunte: l'enunciato della parte è dimostrato
integralmente, nella generalità richiesta (vedi la prova per le ipotesi
esatte). Nessun calcolo al computer è necessario.


% ===== CAMPO: 2. Dimostrazione =====
\paragraph{\texorpdfstring{Cell 1: three pairwise disjoint classes have
a pair of moduli with
\(\gcd\ge3\)}{Cell 1: three pairwise disjoint classes have a pair of moduli with \textbackslash gcd\textbackslash ge3}}\label{cell-1-three-pairwise-disjoint-classes-have-a-pair-of-moduli-with-gcdge3}

\textbf{Target.} If \(a_1\ (\mathrm{mod}\ m_1)\),
\(a_2\ (\mathrm{mod}\ m_2)\), \(a_3\ (\mathrm{mod}\ m_3)\) are pairwise
disjoint congruence classes (\(m_i\ge1\)), then \(\gcd(m_i,m_j)\ge3\)
for some \(i<j\).

\subparagraph{\texorpdfstring{Lemma 0 (the direction of \((*)\) that we
use)}{Lemma 0 (the direction of (*) that we use)}}\label{lemma-0-the-direction-of-that-we-use}

Let \(g=\gcd(m,m')\). If \(g\mid a-a'\) then the classes
\(a\ (\mathrm{mod}\ m)\) and \(a'\ (\mathrm{mod}\ m')\) have a common
element. Equivalently: if the two classes are disjoint then
\(g\nmid a-a'\). \emph{Proof.} Write \(m=g\mu\), \(m'=g\mu'\) with
\(\gcd(\mu,\mu')=1\), and \(a-a'=g\delta\) with \(\delta\in\mathbb Z\).
An element of the first class has the form \(a+mt=a+g\mu t\),
\(t\in\mathbb Z\); it lies in the second class iff
\(m'\mid a+g\mu t-a'\), i.e.~\(g\mu'\mid g(\delta+\mu t)\),
i.e.~\(\mu'\mid\delta+\mu t\). Since \(\gcd(\mu,\mu')=1\), \(\mu\) is
invertible modulo \(\mu'\), so \(t\equiv-\delta\mu^{-1}\pmod{\mu'}\)
solves it. \(\square\) (This is the statement \((*)\) quoted in the
problem; we include the short proof so that the argument is
self-contained.)

\subparagraph{Proof of the cell}\label{proof-of-the-cell}

Let the three classes be pairwise disjoint and suppose, for a
contradiction, that \(\gcd(m_i,m_j)\le2\) for all three pairs. 1.
\emph{Every pairwise gcd equals \(2\).} If \(\gcd(m_i,m_j)=1\) then
\(1\mid a_i-a_j\), so by Lemma 0 the classes \(i\) and \(j\) meet,
contradicting disjointness. Hence \(\gcd(m_i,m_j)=2\) for all \(i<j\).
2. \emph{The residues are pairwise of different parity.} For \(i<j\),
disjointness and Lemma 0 give \(\gcd(m_i,m_j)=2\nmid a_i-a_j\),
i.e.~\(a_i\not\equiv a_j\pmod 2\). 3. \emph{Contradiction.}
\(a_1,a_2,a_3\) would be three integers with pairwise different residues
modulo \(2\); but there are only two residues modulo \(2\), so two of
them coincide (pigeonhole).

Hence the assumption fails: some pair satisfies \(\gcd(m_i,m_j)\ge3\).
\(\blacksquare\)

\emph{Remark.} Step 1 is the first ``trivial bound'' of the problem
statement; it is re-derived here only because the argument needs the
exact value \(2\) of every gcd, not merely \(\ge2\).


% ===== CAMPO: 3. Verifica: istruzioni, dipendenze, tempi =====
Prova a mano, nessuna dipendenza. Verifica meccanica non necessaria; il
Referee automatico (due giudici indipendenti,
\texttt{runs/p4\_c1/attempts/packet\_001.json}) non ha trovato passaggi
non giustificati.


% ===== CAMPO: 4. Fonti e contributo =====
\begin{itemize}
\tightlist
\item
  arXiv:2607.24655 Fornal, Sun, On the problem of large gcd for disjoint
  residue classes (context only: asymptotic result, not used)
\item
  Posizione rispetto allo stato dell'arte: Fornal-Sun {[}2607.24655{]}
  prove an asymptotic lower bound for the largest gcd; the small cases
  k=3,4 are elementary and are not treated there. Our proof is
  independent.
\item
  Contributo: la prova qui scritta è del team, per esteso.
\end{itemize}


% ===== CAMPO: 5. Limiti e parti irrisolte =====
Nessuna: la parte è dimostrata. Gap dichiarati: nessuno.

